\de{ĐỀ THI HỌC KỲ II NĂM HỌC 2022-2023}{THPT Phú Hòa}


\begin{bt}%[0D3B2-3]%[Dự án đề kiểm tra HKII NH22-23- Phạm Văn Long]%[THPT Phú Hoà]
	\begin{enumerate}
		\item Cho hàm số $y=f(x)$ có đồ thị như hình vẽ. Lập bảng xét dấu của $f(x)$.
		\begin{center}
			\begin{tikzpicture}[scale=0.6, font=\footnotesize,line join=round, line cap=round, >=stealth];
				\draw[->] (-4,0) -- (2.5,0) node[below] {\scriptsize $x$};
				\draw[->] (0,-2) -- (0,5) node[left] {\scriptsize $y$};
				\draw (0.1,0.1)node[below left]{\scriptsize $O$};
				\draw[thin,samples=150,smooth,domain=-3.3:1.3] plot(\x,{-1*((\x)^2)-2*(\x)+3});						
				\draw[fill=black]  (-3,0) circle (1.2pt) (1,0) circle (1.2pt);
				\draw (-3.5,0)node[above]{$-3$} (1.3,0)node[above]{$1$};	
			\end{tikzpicture}
		\end{center}
		\item Giải bất phương trình $x^2-8 x+16 \leq 0$.
	\end{enumerate}
\loigiai{
\begin{enumerate}
	\item Bảng xét dấu của $f(x)$.
			\begin{center}
			\begin{tikzpicture}
				\tkzTabInit[lgt=1.2,espcl=2.5,deltacl=0.6]
				{$x$ /0.6,$y$ /0.6}
				{$-\infty$,$-3$,$1$,$+\infty$}
				\tkzTabLine{,-,$0$,+,$0$,-,}				
			\end{tikzpicture}
		\end{center}
	Từ bảng xét dấu của $f(x)$ ta có
	\begin{itemize}
		\item $f(x)> 0$ khi $x\in \left(-3;1\right)$.
		\item $f(x)<0$ khi $x\in (-\infty;-3)\cup (1;+\infty)$.
		\item $f(x)= 0$ khi $x=-3$ hoặc $x=1$.
	\end{itemize}
	\item Ta có $x^2-8 x+16=0\Leftrightarrow x=4$.\\
	Bảng xét dấu của $f(x)=x^2-8 x+16$.
	\begin{center}
		\begin{tikzpicture}
			\tkzTabInit[lgt=1.2,espcl=2.5,deltacl=0.6]
			{$x$ /0.6,$y$ /0.6}
			{$-\infty$,$4$,$+\infty$}
			\tkzTabLine{,+,$0$,+,}				
		\end{tikzpicture}
	\end{center}
	Từ bảng xét dấu của $f(x)$ suy ra $x^2-8 x+16 \geq 0\, \forall  \in \mathbb{R}$.\\
	Khi đó suy ra $x^2-8 x+16 \leq 0$ khi $x=4$.\\
	Vậy $S=\{4\}$.
\end{enumerate}
}
\end{bt}
\begin{bt}%[0D7B3-2]%[Dự án đề kiểm tra HKII NH22-23- Phạm Văn Long]%[THPT Phú Hoà]
	Giải phương trình $\sqrt{x^2+3x+1}=\sqrt{3x^2+6x-4}$.
	\loigiai{
		\allowdisplaybreaks
		\begin{eqnarray*}
			&&\sqrt{x^2+3x+1}=\sqrt{3x^2+6 x-4}\\
			&\Rightarrow& x^2+3x+1=3x^2+6x-4\\
			&\Rightarrow& 2x^2+3x-5=0\\
			&\Rightarrow& \hoac{&x=1\\&x=-\dfrac{5}{2}.}
		\end{eqnarray*}
	Thử lại ta thấy $x=1$ thỏa mãn và $x=-\dfrac{5}{2}$ không thỏa mãn.\\
	Vậy $S=\{1\}$.
}
\end{bt}
\begin{bt}%[0D8B3-1]%[Dự án đề kiểm tra HKII NH22-23- Phạm Văn Long]%[THPT Phú Hoà]
	Khai triến biều thức $(3x-4)^5$.
	\loigiai{
		\allowdisplaybreaks
		\begin{eqnarray*}
			(3x-4)^5&=&(3x)^5+5\cdot (3x)^4\cdot (-4)+10\cdot (3x)^3\cdot (-4)^2+10\cdot (3x)^2\cdot (-4)^3\\
			&&+5\cdot (3x)\cdot (-4)^4+(-4)^5\\
			&=& 243x^5-1620x^4+4320x^3-5760x^2+3840x-1024.
		\end{eqnarray*}	
}
\end{bt}
\begin{bt}%[0D8B2-1]%[Dự án đề kiểm tra HKII NH22-23- Phạm Văn Long]%[THPT Phú Hoà]
	Cho tập hợp $X=\{1 ; 2 ; 3 ; 4 ; 5\}$.
	\begin{enumerate}
		\item Có bao nhiêu tập con gồm ba phần từ cùa $X$?
		\item Có thể lập được bao nhiêu số tự nhiên có 4 chữ số khác nhau từ tập $X$?
	\end{enumerate}
	\loigiai{
\begin{enumerate}
	\item Số tập con gồm ba phần tử cùa $X$ là $\mathrm{C}_5^3=10$.
	\item Số tự nhiên có 4 chữ số khác nhau từ tập $X$ là $\mathrm{A}_5^4=60$ số.
\end{enumerate}	
}
\end{bt}
\begin{bt}%[0D0B2-2]%[Dự án đề kiểm tra HKII NH22-23- Phạm Văn Long]%[THPT Phú Hoà]
	\begin{enumerate}
		\item Tung một đồng xu cân đối và đồng chất 3 lần liên tiếp. Hãy xác định không gian mẫu.
		\item Gieo hai con xúc xắc cân đối và đồng chất. Tính xác suất của biến cố $A$: \lq\lq Tổng số chấm trên hai mặt xuất hiện bằng 5\rq\rq.
	\end{enumerate}
	\loigiai{
\begin{enumerate}
	\item $\Omega=\{SSS;SSN;SNS;NSS;NNS;NSN;SNN;NNN\}$.
	\item $n(\Omega)=36$.\\
	Ta có $A=\{(1;4);(2;3);(3;2);(4;1)\}\Rightarrow n(A)=4$.\\
	Vậy $P(A)=\dfrac{4}{36}=\dfrac{1}{9}$.
\end{enumerate}	
}
\end{bt}

	\begin{bt}%[0T8K2-2]%[0T0B2-2]%[Dự án đề kiểm tra HKII NH22-23- Dương Phước Sang]%[THPT Phú Hoà]
		\indent
		\begin{listEX}
			\item Có bao nhiêu cách chia hết $4$ đồ vật khác nhau cho $3$ người, biết rằng mỗi người nhận ít nhất một đồ vật.
			\item An, Bình và $4$ bạn nữa xếp ngẫu nhiên thành một hàng ngang. Tính xác suất để An và Bình đứng ở hai đầu hàng.
		\end{listEX}
		\loigiai{
			\indent
			\begin{listEX}
				\item Chia hết $4$ đồ vật khác nhau cho $3$ người (mỗi người ít nhất một đồ vật), nhất định có $1$ người nhận được $2$ món đồ đồng thời $2$ người còn lại, mỗi người chỉ nhận được $1$ món.\\
				Để thực hiện được công việc chia đồ đó, trước hết, ta chọn $2$ trong $4$ món đồ vật, rồi chọn $1$ trong $3$ người để nhận nó, có $\mathrm{C}_4^2\cdot \mathrm{C_3^1}$ cách thực hiện.\\
				Còn lại $2$ món đồ vật, chỉ cần chia mỗi món cho $1$ người chưa nhận đồ là được, có $2!$ cách chia đồ vật.\\
				Vậy có $\mathrm{C}_4^2\cdot \mathrm{C_3^1}\cdot 2!=36$ cách chia đồ vật thỏa mãn yêu cầu của bài toán.
				\item Phép thử ``\textit{xếp ngẫu nhiên An, Bình và $4$ bạn nữa thành một hàng ngang}'' có số phần tử của không gian mẫu là $n(\Omega)=6!$.\\
				Gọi $A$ là biến cố ``\textit{An và Bình đứng ở hai đầu hàng}''.\\
				Khi $A$ xảy ra, An và Bình được xếp ở hai đầu hàng đồng thời $4$ bạn còn lại đứng ở $4$ vị trí còn lại. Số kết quả thuận lợi cho $A$ là $n(A)=2!\cdot 4!$.\\
				Vậy xác suất của biến cố $A$ là $\mathrm{P}(A)=\dfrac{n(A)}{n(\Omega)}=\dfrac{2!\cdot 4!}{6!}=\dfrac{1}{15}$.
			\end{listEX}
		}
	\end{bt}
	
	\begin{bt}%[0T9Y2-5]%[0T9B2-2]%[Dự án đề kiểm tra HKII NH22-23- Dương Phước Sang]%[THPT Phú Hoà]
		Cho điểm $M(-1;2)$ và đường thẳng $\Delta\colon 4x-3y-5=0$.
		\begin{listEX}
			\item Tính khoảng cách từ điểm $M$ đến đường thẳng $\Delta$.
			\item Lập phương trình tham số của đường thẳng $d$ đi qua điểm $M$ và vuông góc với $\Delta$.
		\end{listEX}
		\loigiai{
			\indent
			Xét điểm $M(-1;2)$ và đường thẳng $\Delta\colon 4x-3y-5=0$.
			\begin{listEX}
				\item Ta có $\mathrm{d}(M,\Delta)=\dfrac{\big|4(-1)-3\cdot 2-5\big|}{\sqrt{4^2+(-3)^2}}=3$.
				\item Gọi $d$ là đường thẳng $d$ đi qua điểm $M(-1;2)$ và vuông góc với $\Delta\colon 4x-3y-5=0$.\\
				Khi đó $d$ nhận véc-tơ pháp tuyến $\vec{n}_{\Delta}$ của $\Delta$ làm một véc-tơ chỉ phương.\\
				Do đó $d$ có một véc-tơ chỉ phương $\vec{u}_d=\vec{n}_{\Delta}=(4;-3)$.\\
				Ngoài ra $d$ đi qua điểm $M$ nên phương trình tham số của $d\colon \heva{&x=-1+4t\\&y=2-3t.}$
			\end{listEX}
		}
	\end{bt}
	
	\begin{bt}%[0T9Y3-1]%[0T9B3-2]%[Dự án đề kiểm tra HKII NH22-23- Dương Phước Sang]%[THPT Phú Hoà]
		\indent
		\begin{listEX}
			\item Tìm tọa độ tâm $I$ và tính bán kính $R$ của đường tròn $(C)\colon (x-3)^2+(y+4)^2=25$.
			\item Viết phương trình đường tròn $(C')$ có tâm $K(1;-2)$ và đi qua gốc tọa độ $O$.
		\end{listEX}
		\loigiai{
			\indent
			\begin{listEX}
				\item Đường tròn $(C)\colon (x-3)^2+(y+4)^2=25$ có tâm $I(3;-4)$ và bán kính $R=\sqrt{25}=5$.
				\item Đường tròn $(C')$ có tâm $K(1;-2)$ và đi qua $O(0;0)$ có bán kính $R=OK=\sqrt{1^2+(-2)^2}=\sqrt{5}$.\\
				Do đó phương trình đường tròn $(C')\colon (x-1)^2+(y+2)^2=5$.
			\end{listEX}
		}
	\end{bt}
	
	\begin{bt}%[0T9B2-2]%[0T9B3-3]%[Dự án đề kiểm tra HKII NH22-23- Dương Phước Sang]%[THPT Phú Hoà]
		\indent
		\begin{listEX}
			\item Cho tam giác $ABC$, biết $A(-3;0)$, $B(-2;6)$, $C(4;-2)$. Lập phương trình tổng quát của đường trung tuyến $AN$.
			\item Viết phương trình tiếp tuyến của đường tròn $(T)\colon x^2+y^2-2x-4y=0$ tại điểm $E(-1;3)$.
		\end{listEX}
		\loigiai{
			\indent
			Xét $A(-3;0)$, $B(-2;6)$, $C(4;-2)$.
			\begin{listEX}
				\item Trung tuyến $AN$ đi qua trung điểm $N(1;2)$ của cạnh $BC$.\\
				Suy ra $AN$ nhận $\vec{AN}=(4;2)$ và $\vec{u}=(2;1)$ làm các véc-tơ chỉ phương.\\
				Từ đó $AN$ nhận $\vec{n}=(1;-2)$ làm véc-tơ pháp tuyến.\\
				Phương trình tổng quát của $AN\colon 1(x+3)-2(y-0)=0 
				\Leftrightarrow x-2y+3=0$.
				\item Đường tròn $(T)\colon x^2+y^2-2x-4y=0$ có tâm $I(1;2)$ và điểm $E(-1;3) \in (T)$.\\
				Tiếp tuyến của $(T)$ tại $E$ có phương trình
				$$(1+1)(x+1)+(2-3)(y-3)=0 
				\Leftrightarrow 2x-y+5=0.$$
			\end{listEX}
		}
	\end{bt}
	
	\begin{bt}%[0T0K2-2]%[0T8K2-6]%[Dự án đề kiểm tra HKII NH22-23- Dương Phước Sang]%[THPT Phú Hoà]
		\indent
		\begin{listEX}
			\item Một hộp chứa $20$ tấm thẻ có kích thước như nhau và được đánh số từ $1$ đến $20$. Chọn ngẫu nhiên đồng thời $4$ thẻ từ hộp. Tính xác suất để tích các số ghi trên các thẻ chia hết cho $3$.
			\item Một đa giác đều có $2n$ cạnh nội tiếp đường tròn $(O)$. Biết số tam giác tạo thành từ các đỉnh nhiều gấp $20$ lần số hình chữ nhật tạo thành từ các đỉnh. Hỏi đa giác đều đó có tất cả bao nhiêu cạnh?
		\end{listEX}
		\loigiai{
			\indent
			\begin{listEX}
				\item Xét phép thử ``\textit{chọn ngẫu nhiên đồng thời $4$ thẻ từ hộp $20$ thẻ có đánh số từ $1$ đến $20$}''.\\
				Số kết quả có thể xảy ra của phép thử trên là $n(\Omega)=\mathrm{C}_{20}^4$.\\
				Xét biến cố $A$: ``\textit{tích các số ghi trên các thẻ được chọn là một số chia hết cho $3$}''.\\
				$A$ không xảy ra khi ta chọn được cả $4$ thẻ có số trên đó đều không chia hết cho $3$, tức là $4$ số chọn được cùng thuộc tập hợp $\big\{x \in \mathbb{N}^* \mid x \leq 20\big\} \setminus \{3k \mid k \in \mathbb{N}^* \text{ và } k \leq 6\}$.\\
				Do đó, số kết quả thuận lợi cho $A$ là $n(A)=\mathrm{C}_{20}^4-\mathrm{C}_{14}^4$.\\
				Vậy xác suất của biến cố $A$ là $\mathrm{P}(A)=\dfrac{n(A)}{n(\Omega)}=\dfrac{\mathrm{C}_{20}^4-\mathrm{C}_{14}^4}{\mathrm{C}_{20}^4}=\dfrac{3844}{4845}$.
				\item Đường tròn $(O)$ ngoại tiếp đa giác đều $(\mathcal{H})$ $2n$ cạnh có $n$ đường kính đi qua đỉnh của $(\mathcal{H})$.\\
				Mỗi hình chữ nhật lấy đỉnh từ đỉnh của $(\mathcal{H})$ có $2$ đường chéo là $2$ trong $n$ đường kính nêu trên, do đó số hình chữ nhật lấy đỉnh từ đỉnh của $(\mathcal{H})$ là $\mathrm{C}_n^2$.\\
				Cũng từ đỉnh của $(\mathcal{H})$, lập ra các tam giác ta được số tam giác là $\mathrm{C}_{2n}^3$.\\
				Giả thiết bài toán suy ra $\mathrm{C}_{2n}^3=20\mathrm{C}_n^2$ $(n \in \mathbb{N},n \geq 3)$.\\
				Ta có 
				$\begin{aligned}[t]
					\mathrm{C}_{2n}^3=20\mathrm{C}_n^2
					&\Leftrightarrow \dfrac{2n(2n-1)(2n-2)}{3!}=20\cdot\dfrac{n(n-1)}{2!}\\
					&\Leftrightarrow \dfrac{2n(2n-1)(2n-2)}{n(n-1)}=20\cdot\dfrac{3!}{2!}\\
					&\Leftrightarrow 8n-4=60 
					\Leftrightarrow 8n=64
					\Leftrightarrow n=8 \text{ (nhận).}
				\end{aligned}$\\
				Vậy đa giác đều đã cho có $2n=16$ cạnh.
			\end{listEX}
		}
	\end{bt}